% BÀI TOÁN HÌNH HỌC - BAI_HINH_15 (TRÍCH ĐỀ 15)
\begin{ex}
% Câu 14 (Hình học)
	Cho $\triangle PQR$ nhọn, đường tròn tâm $O$ có đường kính $QR$ cắt $PQ, PR$ lần lượt ở $D$ và $E$. Gọi $H$ là giao điểm của $QE$ và $DR$.
	\begin{enumerate}
		\item[1)] Tính số đo $\widehat{QDR}$ và $\widehat{QER}$.
		\item[2)] 
		\begin{enumerate}
			\item[a)] Gọi $I$ là trung điểm của $PH$. Chứng minh: $ID$ là tiếp tuyến của đường tròn $(O)$.
			\item[b)] Gọi $M$ là trung điểm của $HR$. Chứng minh: $IM \perp OM$.
		\end{enumerate}
	\end{enumerate}
\end{ex}
